\documentclass[11pt]{article}

% One-page research showcase for AutoLabs Experiment 3C.
% Title is a placeholder pending the owner's pick (candidates listed at the foot).
% \pending{} values are confirmed by the full-power screen (run …ded5477e).

\usepackage[margin=0.9in]{geometry}
\usepackage{amsmath,amssymb}
\usepackage{xcolor}
\usepackage{tcolorbox}
\usepackage[hidelinks]{hyperref}
\usepackage{titlesec}
\usepackage{enumitem}
\setlist{nosep,leftmargin=1.4em}
\pagestyle{empty}
\newcommand{\pending}[1]{\textcolor{gray}{#1}}
\definecolor{ink}{HTML}{111312}
\definecolor{rust}{HTML}{A85238}

\begin{document}
\begin{center}
{\LARGE\bfseries Unsupervised Discovery of Persona-Relevant Directions\\[2pt]
with Matryoshka Sparse Autoencoders}\\[6pt]
{\normalsize Raphael Khalid \quad$\cdot$\quad AutoLabs \quad$\cdot$\quad
\texttt{autolabs-ebon.vercel.app}}\\[2pt]
{\small\itshape Working draft, September 2026}
\end{center}

\vspace{4pt}
\hrule
\vspace{8pt}

\noindent\textbf{Abstract.}\quad
Persona vectors---activation directions that control traits such as sycophancy or
dishonesty---are found by \emph{naming} a trait, writing contrastive prompts, and
differencing activations, which requires the trait to be specified in advance,
precisely described, and prompt-inducible. Chen et al.\ (2025) conjecture that
``SAEs may therefore enable unsupervised discovery of persona-relevant directions,
including specific traits that cannot be easily elicited through prompting.'' We
take up that conjecture at scale: we train a Matryoshka BatchTopK sparse
autoencoder ($235$M parameters, $32{,}768$ features) on $100$M assistant-position
tokens of \texttt{Qwen2.5-7B-Instruct} at layer~$19$ (held-out FVE~$0.72$,
$L_0{=}40$), rank its dictionary for persona-relevance with three label-free arms,
steer the model along each candidate, and screen the steered behaviour against
random-direction nulls and persona-vector \emph{controls} before a blinded judge
names the survivors---no trait specified and no description written at any point.
A first budget-constrained screen surfaced $\pending{96}$ candidate directions that
separate personas across scenarios (residual-AUC up to $1.0$)---e.g.\ a
poetic/effusive voice and a cautious-analyst voice. A powered confirmation screen
\pending{[in progress]} certifies which clear the control gate and names them. We do
\emph{not} claim these directions are unreachable by prompting; measuring the
\emph{graded} form of Chen et al.'s ``easily''---how hard a prompt must work to
reproduce each---is future work.

\vspace{10pt}
\begin{tcolorbox}[colback=black!3,colframe=ink,boxrule=0.5pt,left=10pt,right=10pt,top=8pt,bottom=8pt]
\textbf{At a glance}
\begin{itemize}
  \item[$\bullet$] \textbf{Dictionary:} $235$M-param Matryoshka BatchTopK SAE, width $32{,}768$, $k{=}40$, nested $1$k/$4$k/$16$k/$32$k shells.
  \item[$\bullet$] \textbf{Trained on:} $100$M assistant-position tokens of Qwen2.5-7B-Instruct, layer $19$ $\rightarrow$ \textbf{held-out FVE $0.72$, $L_0=40$}.
  \item[$\bullet$] \textbf{Discovery:} three \emph{label-free} ranking arms; no trait ever specified.
  \item[$\bullet$] \textbf{Screen:} steering + separability AUC vs $20$ random nulls, gated by persona-vector controls; blinded LLM judge names survivors.
  \item[$\bullet$] \textbf{Drops requirements 1--2} of prompt-based extraction: no trait specified, no description written.
  \item[$\bullet$] \textbf{Open question (future work):} \emph{how easily} a prompt can reproduce each discovered direction---a graded elicitation cost, not a ``prompt-unreachable'' claim.
  \item[$\bullet$] \textbf{Pilot signal:} $\pending{484}$/$512$ direction-signs beat the null ceiling; $\pending{96}$/$256$ features cleared the full gate.
\end{itemize}
\end{tcolorbox}

\vspace{6pt}
\noindent\footnotesize
\textbf{Infrastructure.} Run end-to-end on a reproducible harness (durable
workflows, frozen protocol hashes, a spend ledger and cost caps); the SAE is
trained once on an A100 and reused off a Hugging Face checkpoint for every screen.
\quad\textbf{Builds on} R.~Chen, A.~Arditi, H.~Sleight et al., \emph{Persona
Vectors: Monitoring and Controlling Character Traits in Language Models},
arXiv:2507.21509 (2025), whose appendix first suggests this SAE direction.

\end{document}
